\subsection{Correlated traces in biquadratic dimension fields}
Let $K=\mathbb Q(\sqrt q,\sqrt r)$, with coprime squarefree integers
$q,r>1$. For $s\in\{q,r,qr\}$ put $F_s=s$ if $s\equiv1\pmod4$,
and $F_s=4s$ otherwise, and let $\chi_s$ be the corresponding primitive
quadratic character. Coordinates below use
$(1,\sqrt q,\sqrt r,\sqrt{qr})$.

\begin{proposition}[Complete finite-height relative traces]
\label{prop:m3-biquadratic-traces}
The normalized relative trace of a primitive $n$th root to $K$ has rational
coordinate $\mu(n)/\varphi(n)$. Its $\sqrt s$ coordinate is zero unless
$F_s\mid n$ and $\gcd(n/F_s,F_s)=1$; otherwise it equals
\[
 \frac{\epsilon_s\mu(n/F_s)\chi_s(n/F_s)\sqrt{F_s/s}}
      {\varphi(n)}.
\]
There are four linked sign choices, with
$\epsilon_{qr}=\epsilon_q\epsilon_r$.
For the $k$th power of a primitive $N$th root, set $g=\gcd(N,k)$ and
$n=N/g$, use the same fixed sign pair, and multiply the quadratic
coordinate by $\chi_s(k/g)$.

Suppose the prime support of twist orders is contained in $P$.
For a height $h$, put
\[
 B_h=\prod_{p\in P}
 p^{\max(1,v_p(F_q),v_p(F_r),v_p(F_{qr}))+\lfloor\log_p h\rfloor}.
\]
Every joint relative trace of powers $1,\ldots,h$ is represented by an
order dividing $B_h$, with one of the four linked sign choices, or by
the joint zero vector.
\end{proposition}
\begin{proof}
The rational projection is the Ramanujan sum $\mu(n)$. The quadratic
projections are the primitive quadratic Gauss sums
$\tau(\chi_s)=\sqrt{F_s}$ followed by the coprime-part Ramanujan sum.
Dividing by $\varphi(n)\sqrt s$ gives the displayed formula.
The three signs come from one field automorphism and therefore satisfy
the stated product relation. Reduction of the root order gives the power
rule; independent signs for different powers would discard essential
correlations.

If $v_p(N)$ exceeds the exponent in $B_h$, then for every $k\leq h$
the $p$-valuation of $N/\gcd(N,k)$ exceeds one and every conductor
valuation. The rational M\"obius term vanishes. Each quadratic term
vanishes either by noncoprimality with its conductor or by a zero
M\"obius value. All retained powers consequently have zero trace.
Every other order divides $B_h$.
\end{proof}
For a pseudo-unitary category the spherical Cauchy theorem
\cite{BNRWCauchy} supplies $P$ from the norm of the global dimension, after
the exact integrality and conductor-support checks. These are bounds on
retained traces, not on the unknown categorical exponent.
The field $\mathbb Q(\sqrt2,\sqrt3)$ at height six gives $B_h=288$;
$(q,r)=(3,7),(3,5),(3,13)$ at height four give respectively
$1008,720,1872$ for the prime supports recorded in the certificates.

\paragraph{Weighted induction and exact separation.}
The full induction Gram matrix is computed from the tensor, including
noncommutative rings. For each simple Grothendieck-algebra representation
$E$, its distinguished central constituent has dimension $D/f_E$ and
multiplicity $\dim E$ in $I(1)$; the corresponding restriction column
starts with $\dim E$. Its outer square contributes once to the Gram
matrix. Thus a two-dimensional representation is handled with first
entry two, rather than a commutative first-entry-one shortcut.
All compatible low columns and all residual Gram-bounded nonnegative
columns are enumerated. Their dimensions satisfy the necessary
integrality and $d$-number tests.

Nonnegative weights for high columns and the complete linked trace
types obey the Gram equations, indicator identities and cyclic-rotation
traces on $\operatorname{Hom}(\mathbf1,X^m)$ through the chosen height.
Every pseudo-unitary categorification supplies a feasible point. An
exact rational Farkas witness therefore proves a contradiction; the
numerical solver only proposes the witness.

\paragraph{General scope and independent checks.}
Such a rotation contradiction by itself has pseudo-unitary scope.
General exclusion additionally checks every original dimension
character and every oriented signed quotient in sphericalization,
using exact character, parity and gauge-equivalence certificates.
The accepted applications include $\sha{R8_3_347}$,
$\sha{R8_3_916}$ and $\sha{R8_3_1348}$.
The independent checker constructs the cyclotomic quotient
$\mathbb Q[x]/\Phi_{B_h}$ and its quadratic Gauss sums, checks their
squares and correlated product, and compares every bounded-order
trace and retained power with the formula. This verifies the finite
identities; the preceding argument proves completeness beyond the box.
